% Insertion-ready section for the joint prime-bias manuscript.
% The host manuscript should define \spectralfigurepath before % Insertion-ready section for the joint prime-bias manuscript.
% The host manuscript should define \spectralfigurepath before % Insertion-ready section for the joint prime-bias manuscript.
% The host manuscript should define \spectralfigurepath before % Insertion-ready section for the joint prime-bias manuscript.
% The host manuscript should define \spectralfigurepath before \input{section.tex}.
% Example:
% \newcommand{\spectralfigurepath}{path/to/spectral_reconstruction.pdf}

\section{Transient spectral oscillations at the 300-trillion scale}
\label{sec:spectral-transients}

We compare the ordinary prime-count races with low-zero truncations of their
standard explicit formulas.  This calculation has two deliberately limited
purposes.  First, it tests whether low Dirichlet zeros account quantitatively
for the trajectories observed through $3\times 10^{14}$.  Second, it tests the
specific finite-scale claim that the complex-character interference has
subsided by that range.  The calculation concerns the unregularized counting
functions $\pi(x;N,a)$.  It neither proves the asymptotic proposed for the
Gaussian-regularized statistic elsewhere in this paper nor transfers an
ordering between the two statistics.

\subsection{Normalization and spectral truncation}

For a reduced residue class $a\pmod N$, set
\begin{equation}
  E_N(x;a,1)
  :=\frac{\varphi(N)\log x}{\sqrt{x}}
    \bigl(\pi(x;N,a)-\pi(x;N,1)\bigr).
  \label{eq:ordinary-race-normalization}
\end{equation}
Let
\[
  C_N(a):=\#\{b\bmod N:b^2\equiv a\pmod N\}.
\]
Assuming the critical-line form used in the standard GRH explicit formula,
we form the finite truncation
\begin{equation}
  \widetilde E_{N,K}(x;a,1)
  :=C_N(1)-C_N(a)
   -\sum_{\chi\ne\chi_0}\sum_{j=1}^{K}
     2\operatorname{Re}\!\left(
       \frac{(\overline{\chi(a)}-1)e^{i\gamma_{\chi,j}\log x}}
            {\tfrac12+i\gamma_{\chi,j}}
     \right),
  \label{eq:ordinary-race-truncation}
\end{equation}
where $\gamma_{\chi,j}>0$ is the $j$th positive ordinate returned for
$L(s,\chi)$.  Thus $K$ means the first $K$ positive ordinates for
\emph{each} nonprincipal character, rather than the first $K$ modes globally.
The coefficient $\overline{\chi(a)}-1$ is consistent with selecting the
ordered difference between the classes $a$ and $1$.  It should not be
confused with the separate character kernel used to define the regularized
statistic.

The ordinates were computed for every nonprincipal character modulo
\[
  N\in\{7,8,11,19,23\}
\]
through height $80$.  Two runs of \texttt{lfunzeros}, with mesh parameters
$64$ and $96$, were required to agree character by character.  Every returned
ordinate used in the reconstruction was also required to satisfy
\[
  \left|L\!\left(\tfrac12+i\gamma_{\chi,j},\chi\right)\right|<10^{-28}.
\]
The two mesh computations are robustness checks within one implementation,
not an independent proof that the zero list is complete.  As a separate
cross-check, all three first-zero anchors modulo $8$ agree with a pre-existing
high-precision \texttt{mpmath} computation to $10^{-12}$.

For the exceptional modulus-$19$ mode discussed below, a second implementation
constructs the character directly with Python FLINT and evaluates its Hardy
$Z$-function using Arb ball arithmetic.  Opposite, sign-definite endpoint
balls give an independent existence certificate for a critical-line zero.
This certificate does not assert uniqueness of that zero, completeness below
it, or GRH.

\subsection{Finite computational result}

The observed data consist of 438 logarithmically spaced checkpoints through
$3\times10^{14}$.  The ``top decade'' below is the 53-point subwindow
$x\ge 3\times10^{13}$.

\begin{compresult}[verified finite computation]
\label{prop:finite-spectral-computation}
For the immutable input dataset identified in the reproducibility statement
below, the checked pipeline produces 13,578 reconstruction rows.  At
$K=25$, the correlations between
$\widetilde E_{N,25}(x;-1,1)$ and $E_N(x;-1,1)$ on the 53-point top-decade
window are respectively
\[
  0.964909993,\quad 0.930923926,\quad 0.938516494,
  \quad 0.971455692,\quad 0.826435845
\]
for $N=7,8,11,19,23$.  Across the 52 consecutive top-decade intervals,
the observed rank of $-1$ changes respectively
\[
  17,\quad 23,\quad 30,\quad 17,\quad 39
\]
times.  In particular, for every one of these five moduli the class $-1$
both attains rank one and leaves rank one within the top decade.
\end{compresult}

This is a computational proposition: it records the output of a specified,
hash-checked finite calculation.  It is not a theorem about omitted zeros,
GRH, or limiting prime-race behavior.

\begin{table}[tb]
  \centering
  \small
  \caption{Top-decade reconstruction of the $-1$ race.  The correlation and
  RMSE use $K=25$ zeros per nonprincipal character.  Rank and leader changes
  refer to the observed ordinary-count curves over 52 consecutive intervals.}
  \label{tab:top-decade-spectral-fit}
  \begin{tabular}{@{}rrrrrr@{}}
    \toprule
    $N$ & correlation & RMSE & rank changes & leader changes
        & fraction $-1$ leads\\
    \midrule
     7 & 0.965 & 0.515 & 17 & 13 & 0.358\\
     8 & 0.931 & 0.626 & 23 & 15 & 0.358\\
    11 & 0.939 & 0.919 & 30 & 26 & 0.377\\
    19 & 0.971 & 1.745 & 17 &  8 & 0.415\\
    23 & 0.826 & 1.728 & 39 & 23 & 0.132\\
    \bottomrule
  \end{tabular}
\end{table}

Figure~\ref{fig:spectral-reconstruction} compares the observed curves with
the $K=1,3,10,25$ truncations.  For $N=7,11,19$, even a few modes reproduce
much of the slow structure.  For $N=23$, one zero per character is inadequate,
and the fit improves only when substantially more modes are retained.  This
behavior makes a single-mode ``onset scale'' unreliable: distinct characters
and ordinates remain active simultaneously.

\begin{figure}[p]
  \centering
  \includegraphics[width=\textwidth]{\spectralfigurepath}
  \caption{Observed normalized ordinary-count races (black), compared with
  explicit-formula truncations using the first $K=1,3,10,25$ positive zeros of
  every nonprincipal character.  The panels show $N=7,8,11,19,23$ and the
  class $a=-1\pmod N$.}
  \label{fig:spectral-reconstruction}
\end{figure}

\subsection{The very low odd-character mode modulo 19}

The computation identifies a load-bearing correction to a one-mode
description of the modulus-19 race.  For the primitive character represented
by Conrey index $13$ modulo $19$, of character order $18$, the first returned
positive ordinate is
\begin{equation}
  \gamma_{19,13,1}=0.0189563990802261\ldots,
  \qquad \chi(-1)=-1.
  \label{eq:mod19-low-zero}
\end{equation}
Both mesh runs return this ordinate, and direct evaluation at the zero passes
the residual threshold above.  Independently, Arb gives a sign-change bracket
of width $5.94\times10^{-84}$ around
$0.018956399080226142994\ldots$ and an upper bound below
$3.55\times10^{-85}$ for the modulus of the midpoint $L$-value.  The
intermediate-value theorem therefore certifies existence of at least one zero
in the bracket.  The corresponding mode is active in the $-1$ race: on the
top decade its root-mean-square contribution is $6.719$, and its correlation
with the centered observed curve is $0.728$.

The conclusion is stable under a substantially deeper truncation.  Extending
all $17$ nonprincipal characters modulo $19$ to at least $100$ positive
ordinates improves the top-decade $-1$ correlation from $0.9715$ at $K=25$
to $0.9925$ at $K=100$, while the RMSE decreases from $1.745$ to $1.252$.
The $K=100$ reconstruction retains $14$ rank changes and $7$ leader changes,
with rank and leader agreement against the observed data of $90.6\%$ and
$98.1\%$, respectively.  Increasing spectral depth therefore strengthens,
rather than removes, the transient-rank conclusion.

Consequently, an estimate based on an ordinate near $1.74$ cannot be described
as using the lowest relevant complex-character zero modulo $19$.  The much
smaller ordinate in \eqref{eq:mod19-low-zero} produces a substantially slower
oscillation that remains visible near $3\times10^{14}$.  This does not by
itself settle the ultimate behavior of the race; it shows that the proposed
finite scale is not a demonstrated point beyond complex-character
interference.

\subsection{Observed rankings and what they do not imply}

At the final checkpoint $x=3\times10^{14}$, the class $-1$ ranks first among
the quadratic nonresidues for $N=7$ and $N=23$, third of three for $N=8$,
third of five for $N=11$, and sixth of nine for $N=19$.  Those endpoint ranks
are snapshots, and the transition counts in
Table~\ref{tab:top-decade-spectral-fit} show why they cannot be treated as
stable asymptotic orderings.  The appropriate finite-scale conclusion is:

\begin{quote}
The unregularized races remain spectrally coherent but dynamically unsettled
through $3\times10^{14}$.  Low-zero explicit-formula truncations reproduce the
observed trajectories with high correlation, while frequent rank and leader
changes persist throughout the top decade.  These data support a low-zero
transient mechanism, not a universal stabilization threshold or an eventual
ordinary-count ordering.
\end{quote}

The distinction between statistics is essential.  Equations
\eqref{eq:ordinary-race-normalization}--\eqref{eq:ordinary-race-truncation}
describe ordinary counts.  A theorem about a separately defined
Gaussian-regularized spectral statistic would concern that statistic only.
Inferring an eventual ordering of ordinary counts would require an additional
transfer theorem with explicit error control; no such transfer is supplied by
the computation here.

\subsection{Reproducibility, data availability, and limitations}

The authoritative input is \texttt{curve\_3e14.tsv}, with SHA-256 digest
\begin{center}
\ttfamily\footnotesize
57957bdb3ce3243272c3d4b8e9ffe7dfb734b759f48b63becf7ae6f924e1caab.
\end{center}
The accompanying reproducibility package contains the PARI/GP zero generator,
the reconstruction program, fit metrics, the complete transition ledger,
per-mode attribution, figure source, a negative-test suite for the verifier,
and a SHA-256 manifest.  Its \texttt{independent\_n19} follow-up contains the
Arb sign-change certificate, the $K=25,50,100$ reconstruction, seven
adversarial verifier fixtures, and a separate manifest.  Running the two
pipeline scripts regenerates their tabular outputs and checks the declared
invariants.  For an
archival release, the package, input dataset, exact software versions, and
manifest should be deposited under a permanent DOI; the DOI replaces this
sentence before submission.

The numerical interpretation has four limits.  First, the reconstruction uses
the GRH critical-line form of the explicit formula.  Second, finite zero lists
do not prove GRH or completeness of the zero set.  Third, correlation measures
finite explanatory agreement, not mathematical implication.  Fourth, neither
the observed curves nor their truncations prove a universal or eventual
ordering.  The results are therefore best used as a controlled finite-scale
diagnostic and as a constraint on any proposed asymptotic theory.
.
% Example:
% \newcommand{\spectralfigurepath}{path/to/spectral_reconstruction.pdf}

\section{Transient spectral oscillations at the 300-trillion scale}
\label{sec:spectral-transients}

We compare the ordinary prime-count races with low-zero truncations of their
standard explicit formulas.  This calculation has two deliberately limited
purposes.  First, it tests whether low Dirichlet zeros account quantitatively
for the trajectories observed through $3\times 10^{14}$.  Second, it tests the
specific finite-scale claim that the complex-character interference has
subsided by that range.  The calculation concerns the unregularized counting
functions $\pi(x;N,a)$.  It neither proves the asymptotic proposed for the
Gaussian-regularized statistic elsewhere in this paper nor transfers an
ordering between the two statistics.

\subsection{Normalization and spectral truncation}

For a reduced residue class $a\pmod N$, set
\begin{equation}
  E_N(x;a,1)
  :=\frac{\varphi(N)\log x}{\sqrt{x}}
    \bigl(\pi(x;N,a)-\pi(x;N,1)\bigr).
  \label{eq:ordinary-race-normalization}
\end{equation}
Let
\[
  C_N(a):=\#\{b\bmod N:b^2\equiv a\pmod N\}.
\]
Assuming the critical-line form used in the standard GRH explicit formula,
we form the finite truncation
\begin{equation}
  \widetilde E_{N,K}(x;a,1)
  :=C_N(1)-C_N(a)
   -\sum_{\chi\ne\chi_0}\sum_{j=1}^{K}
     2\operatorname{Re}\!\left(
       \frac{(\overline{\chi(a)}-1)e^{i\gamma_{\chi,j}\log x}}
            {\tfrac12+i\gamma_{\chi,j}}
     \right),
  \label{eq:ordinary-race-truncation}
\end{equation}
where $\gamma_{\chi,j}>0$ is the $j$th positive ordinate returned for
$L(s,\chi)$.  Thus $K$ means the first $K$ positive ordinates for
\emph{each} nonprincipal character, rather than the first $K$ modes globally.
The coefficient $\overline{\chi(a)}-1$ is consistent with selecting the
ordered difference between the classes $a$ and $1$.  It should not be
confused with the separate character kernel used to define the regularized
statistic.

The ordinates were computed for every nonprincipal character modulo
\[
  N\in\{7,8,11,19,23\}
\]
through height $80$.  Two runs of \texttt{lfunzeros}, with mesh parameters
$64$ and $96$, were required to agree character by character.  Every returned
ordinate used in the reconstruction was also required to satisfy
\[
  \left|L\!\left(\tfrac12+i\gamma_{\chi,j},\chi\right)\right|<10^{-28}.
\]
The two mesh computations are robustness checks within one implementation,
not an independent proof that the zero list is complete.  As a separate
cross-check, all three first-zero anchors modulo $8$ agree with a pre-existing
high-precision \texttt{mpmath} computation to $10^{-12}$.

For the exceptional modulus-$19$ mode discussed below, a second implementation
constructs the character directly with Python FLINT and evaluates its Hardy
$Z$-function using Arb ball arithmetic.  Opposite, sign-definite endpoint
balls give an independent existence certificate for a critical-line zero.
This certificate does not assert uniqueness of that zero, completeness below
it, or GRH.

\subsection{Finite computational result}

The observed data consist of 438 logarithmically spaced checkpoints through
$3\times10^{14}$.  The ``top decade'' below is the 53-point subwindow
$x\ge 3\times10^{13}$.

\begin{compresult}[verified finite computation]
\label{prop:finite-spectral-computation}
For the immutable input dataset identified in the reproducibility statement
below, the checked pipeline produces 13,578 reconstruction rows.  At
$K=25$, the correlations between
$\widetilde E_{N,25}(x;-1,1)$ and $E_N(x;-1,1)$ on the 53-point top-decade
window are respectively
\[
  0.964909993,\quad 0.930923926,\quad 0.938516494,
  \quad 0.971455692,\quad 0.826435845
\]
for $N=7,8,11,19,23$.  Across the 52 consecutive top-decade intervals,
the observed rank of $-1$ changes respectively
\[
  17,\quad 23,\quad 30,\quad 17,\quad 39
\]
times.  In particular, for every one of these five moduli the class $-1$
both attains rank one and leaves rank one within the top decade.
\end{compresult}

This is a computational proposition: it records the output of a specified,
hash-checked finite calculation.  It is not a theorem about omitted zeros,
GRH, or limiting prime-race behavior.

\begin{table}[tb]
  \centering
  \small
  \caption{Top-decade reconstruction of the $-1$ race.  The correlation and
  RMSE use $K=25$ zeros per nonprincipal character.  Rank and leader changes
  refer to the observed ordinary-count curves over 52 consecutive intervals.}
  \label{tab:top-decade-spectral-fit}
  \begin{tabular}{@{}rrrrrr@{}}
    \toprule
    $N$ & correlation & RMSE & rank changes & leader changes
        & fraction $-1$ leads\\
    \midrule
     7 & 0.965 & 0.515 & 17 & 13 & 0.358\\
     8 & 0.931 & 0.626 & 23 & 15 & 0.358\\
    11 & 0.939 & 0.919 & 30 & 26 & 0.377\\
    19 & 0.971 & 1.745 & 17 &  8 & 0.415\\
    23 & 0.826 & 1.728 & 39 & 23 & 0.132\\
    \bottomrule
  \end{tabular}
\end{table}

Figure~\ref{fig:spectral-reconstruction} compares the observed curves with
the $K=1,3,10,25$ truncations.  For $N=7,11,19$, even a few modes reproduce
much of the slow structure.  For $N=23$, one zero per character is inadequate,
and the fit improves only when substantially more modes are retained.  This
behavior makes a single-mode ``onset scale'' unreliable: distinct characters
and ordinates remain active simultaneously.

\begin{figure}[p]
  \centering
  \includegraphics[width=\textwidth]{\spectralfigurepath}
  \caption{Observed normalized ordinary-count races (black), compared with
  explicit-formula truncations using the first $K=1,3,10,25$ positive zeros of
  every nonprincipal character.  The panels show $N=7,8,11,19,23$ and the
  class $a=-1\pmod N$.}
  \label{fig:spectral-reconstruction}
\end{figure}

\subsection{The very low odd-character mode modulo 19}

The computation identifies a load-bearing correction to a one-mode
description of the modulus-19 race.  For the primitive character represented
by Conrey index $13$ modulo $19$, of character order $18$, the first returned
positive ordinate is
\begin{equation}
  \gamma_{19,13,1}=0.0189563990802261\ldots,
  \qquad \chi(-1)=-1.
  \label{eq:mod19-low-zero}
\end{equation}
Both mesh runs return this ordinate, and direct evaluation at the zero passes
the residual threshold above.  Independently, Arb gives a sign-change bracket
of width $5.94\times10^{-84}$ around
$0.018956399080226142994\ldots$ and an upper bound below
$3.55\times10^{-85}$ for the modulus of the midpoint $L$-value.  The
intermediate-value theorem therefore certifies existence of at least one zero
in the bracket.  The corresponding mode is active in the $-1$ race: on the
top decade its root-mean-square contribution is $6.719$, and its correlation
with the centered observed curve is $0.728$.

The conclusion is stable under a substantially deeper truncation.  Extending
all $17$ nonprincipal characters modulo $19$ to at least $100$ positive
ordinates improves the top-decade $-1$ correlation from $0.9715$ at $K=25$
to $0.9925$ at $K=100$, while the RMSE decreases from $1.745$ to $1.252$.
The $K=100$ reconstruction retains $14$ rank changes and $7$ leader changes,
with rank and leader agreement against the observed data of $90.6\%$ and
$98.1\%$, respectively.  Increasing spectral depth therefore strengthens,
rather than removes, the transient-rank conclusion.

Consequently, an estimate based on an ordinate near $1.74$ cannot be described
as using the lowest relevant complex-character zero modulo $19$.  The much
smaller ordinate in \eqref{eq:mod19-low-zero} produces a substantially slower
oscillation that remains visible near $3\times10^{14}$.  This does not by
itself settle the ultimate behavior of the race; it shows that the proposed
finite scale is not a demonstrated point beyond complex-character
interference.

\subsection{Observed rankings and what they do not imply}

At the final checkpoint $x=3\times10^{14}$, the class $-1$ ranks first among
the quadratic nonresidues for $N=7$ and $N=23$, third of three for $N=8$,
third of five for $N=11$, and sixth of nine for $N=19$.  Those endpoint ranks
are snapshots, and the transition counts in
Table~\ref{tab:top-decade-spectral-fit} show why they cannot be treated as
stable asymptotic orderings.  The appropriate finite-scale conclusion is:

\begin{quote}
The unregularized races remain spectrally coherent but dynamically unsettled
through $3\times10^{14}$.  Low-zero explicit-formula truncations reproduce the
observed trajectories with high correlation, while frequent rank and leader
changes persist throughout the top decade.  These data support a low-zero
transient mechanism, not a universal stabilization threshold or an eventual
ordinary-count ordering.
\end{quote}

The distinction between statistics is essential.  Equations
\eqref{eq:ordinary-race-normalization}--\eqref{eq:ordinary-race-truncation}
describe ordinary counts.  A theorem about a separately defined
Gaussian-regularized spectral statistic would concern that statistic only.
Inferring an eventual ordering of ordinary counts would require an additional
transfer theorem with explicit error control; no such transfer is supplied by
the computation here.

\subsection{Reproducibility, data availability, and limitations}

The authoritative input is \texttt{curve\_3e14.tsv}, with SHA-256 digest
\begin{center}
\ttfamily\footnotesize
57957bdb3ce3243272c3d4b8e9ffe7dfb734b759f48b63becf7ae6f924e1caab.
\end{center}
The accompanying reproducibility package contains the PARI/GP zero generator,
the reconstruction program, fit metrics, the complete transition ledger,
per-mode attribution, figure source, a negative-test suite for the verifier,
and a SHA-256 manifest.  Its \texttt{independent\_n19} follow-up contains the
Arb sign-change certificate, the $K=25,50,100$ reconstruction, seven
adversarial verifier fixtures, and a separate manifest.  Running the two
pipeline scripts regenerates their tabular outputs and checks the declared
invariants.  For an
archival release, the package, input dataset, exact software versions, and
manifest should be deposited under a permanent DOI; the DOI replaces this
sentence before submission.

The numerical interpretation has four limits.  First, the reconstruction uses
the GRH critical-line form of the explicit formula.  Second, finite zero lists
do not prove GRH or completeness of the zero set.  Third, correlation measures
finite explanatory agreement, not mathematical implication.  Fourth, neither
the observed curves nor their truncations prove a universal or eventual
ordering.  The results are therefore best used as a controlled finite-scale
diagnostic and as a constraint on any proposed asymptotic theory.
.
% Example:
% \newcommand{\spectralfigurepath}{path/to/spectral_reconstruction.pdf}

\section{Transient spectral oscillations at the 300-trillion scale}
\label{sec:spectral-transients}

We compare the ordinary prime-count races with low-zero truncations of their
standard explicit formulas.  This calculation has two deliberately limited
purposes.  First, it tests whether low Dirichlet zeros account quantitatively
for the trajectories observed through $3\times 10^{14}$.  Second, it tests the
specific finite-scale claim that the complex-character interference has
subsided by that range.  The calculation concerns the unregularized counting
functions $\pi(x;N,a)$.  It neither proves the asymptotic proposed for the
Gaussian-regularized statistic elsewhere in this paper nor transfers an
ordering between the two statistics.

\subsection{Normalization and spectral truncation}

For a reduced residue class $a\pmod N$, set
\begin{equation}
  E_N(x;a,1)
  :=\frac{\varphi(N)\log x}{\sqrt{x}}
    \bigl(\pi(x;N,a)-\pi(x;N,1)\bigr).
  \label{eq:ordinary-race-normalization}
\end{equation}
Let
\[
  C_N(a):=\#\{b\bmod N:b^2\equiv a\pmod N\}.
\]
Assuming the critical-line form used in the standard GRH explicit formula,
we form the finite truncation
\begin{equation}
  \widetilde E_{N,K}(x;a,1)
  :=C_N(1)-C_N(a)
   -\sum_{\chi\ne\chi_0}\sum_{j=1}^{K}
     2\operatorname{Re}\!\left(
       \frac{(\overline{\chi(a)}-1)e^{i\gamma_{\chi,j}\log x}}
            {\tfrac12+i\gamma_{\chi,j}}
     \right),
  \label{eq:ordinary-race-truncation}
\end{equation}
where $\gamma_{\chi,j}>0$ is the $j$th positive ordinate returned for
$L(s,\chi)$.  Thus $K$ means the first $K$ positive ordinates for
\emph{each} nonprincipal character, rather than the first $K$ modes globally.
The coefficient $\overline{\chi(a)}-1$ is consistent with selecting the
ordered difference between the classes $a$ and $1$.  It should not be
confused with the separate character kernel used to define the regularized
statistic.

The ordinates were computed for every nonprincipal character modulo
\[
  N\in\{7,8,11,19,23\}
\]
through height $80$.  Two runs of \texttt{lfunzeros}, with mesh parameters
$64$ and $96$, were required to agree character by character.  Every returned
ordinate used in the reconstruction was also required to satisfy
\[
  \left|L\!\left(\tfrac12+i\gamma_{\chi,j},\chi\right)\right|<10^{-28}.
\]
The two mesh computations are robustness checks within one implementation,
not an independent proof that the zero list is complete.  As a separate
cross-check, all three first-zero anchors modulo $8$ agree with a pre-existing
high-precision \texttt{mpmath} computation to $10^{-12}$.

For the exceptional modulus-$19$ mode discussed below, a second implementation
constructs the character directly with Python FLINT and evaluates its Hardy
$Z$-function using Arb ball arithmetic.  Opposite, sign-definite endpoint
balls give an independent existence certificate for a critical-line zero.
This certificate does not assert uniqueness of that zero, completeness below
it, or GRH.

\subsection{Finite computational result}

The observed data consist of 438 logarithmically spaced checkpoints through
$3\times10^{14}$.  The ``top decade'' below is the 53-point subwindow
$x\ge 3\times10^{13}$.

\begin{compresult}[verified finite computation]
\label{prop:finite-spectral-computation}
For the immutable input dataset identified in the reproducibility statement
below, the checked pipeline produces 13,578 reconstruction rows.  At
$K=25$, the correlations between
$\widetilde E_{N,25}(x;-1,1)$ and $E_N(x;-1,1)$ on the 53-point top-decade
window are respectively
\[
  0.964909993,\quad 0.930923926,\quad 0.938516494,
  \quad 0.971455692,\quad 0.826435845
\]
for $N=7,8,11,19,23$.  Across the 52 consecutive top-decade intervals,
the observed rank of $-1$ changes respectively
\[
  17,\quad 23,\quad 30,\quad 17,\quad 39
\]
times.  In particular, for every one of these five moduli the class $-1$
both attains rank one and leaves rank one within the top decade.
\end{compresult}

This is a computational proposition: it records the output of a specified,
hash-checked finite calculation.  It is not a theorem about omitted zeros,
GRH, or limiting prime-race behavior.

\begin{table}[tb]
  \centering
  \small
  \caption{Top-decade reconstruction of the $-1$ race.  The correlation and
  RMSE use $K=25$ zeros per nonprincipal character.  Rank and leader changes
  refer to the observed ordinary-count curves over 52 consecutive intervals.}
  \label{tab:top-decade-spectral-fit}
  \begin{tabular}{@{}rrrrrr@{}}
    \toprule
    $N$ & correlation & RMSE & rank changes & leader changes
        & fraction $-1$ leads\\
    \midrule
     7 & 0.965 & 0.515 & 17 & 13 & 0.358\\
     8 & 0.931 & 0.626 & 23 & 15 & 0.358\\
    11 & 0.939 & 0.919 & 30 & 26 & 0.377\\
    19 & 0.971 & 1.745 & 17 &  8 & 0.415\\
    23 & 0.826 & 1.728 & 39 & 23 & 0.132\\
    \bottomrule
  \end{tabular}
\end{table}

Figure~\ref{fig:spectral-reconstruction} compares the observed curves with
the $K=1,3,10,25$ truncations.  For $N=7,11,19$, even a few modes reproduce
much of the slow structure.  For $N=23$, one zero per character is inadequate,
and the fit improves only when substantially more modes are retained.  This
behavior makes a single-mode ``onset scale'' unreliable: distinct characters
and ordinates remain active simultaneously.

\begin{figure}[p]
  \centering
  \includegraphics[width=\textwidth]{\spectralfigurepath}
  \caption{Observed normalized ordinary-count races (black), compared with
  explicit-formula truncations using the first $K=1,3,10,25$ positive zeros of
  every nonprincipal character.  The panels show $N=7,8,11,19,23$ and the
  class $a=-1\pmod N$.}
  \label{fig:spectral-reconstruction}
\end{figure}

\subsection{The very low odd-character mode modulo 19}

The computation identifies a load-bearing correction to a one-mode
description of the modulus-19 race.  For the primitive character represented
by Conrey index $13$ modulo $19$, of character order $18$, the first returned
positive ordinate is
\begin{equation}
  \gamma_{19,13,1}=0.0189563990802261\ldots,
  \qquad \chi(-1)=-1.
  \label{eq:mod19-low-zero}
\end{equation}
Both mesh runs return this ordinate, and direct evaluation at the zero passes
the residual threshold above.  Independently, Arb gives a sign-change bracket
of width $5.94\times10^{-84}$ around
$0.018956399080226142994\ldots$ and an upper bound below
$3.55\times10^{-85}$ for the modulus of the midpoint $L$-value.  The
intermediate-value theorem therefore certifies existence of at least one zero
in the bracket.  The corresponding mode is active in the $-1$ race: on the
top decade its root-mean-square contribution is $6.719$, and its correlation
with the centered observed curve is $0.728$.

The conclusion is stable under a substantially deeper truncation.  Extending
all $17$ nonprincipal characters modulo $19$ to at least $100$ positive
ordinates improves the top-decade $-1$ correlation from $0.9715$ at $K=25$
to $0.9925$ at $K=100$, while the RMSE decreases from $1.745$ to $1.252$.
The $K=100$ reconstruction retains $14$ rank changes and $7$ leader changes,
with rank and leader agreement against the observed data of $90.6\%$ and
$98.1\%$, respectively.  Increasing spectral depth therefore strengthens,
rather than removes, the transient-rank conclusion.

Consequently, an estimate based on an ordinate near $1.74$ cannot be described
as using the lowest relevant complex-character zero modulo $19$.  The much
smaller ordinate in \eqref{eq:mod19-low-zero} produces a substantially slower
oscillation that remains visible near $3\times10^{14}$.  This does not by
itself settle the ultimate behavior of the race; it shows that the proposed
finite scale is not a demonstrated point beyond complex-character
interference.

\subsection{Observed rankings and what they do not imply}

At the final checkpoint $x=3\times10^{14}$, the class $-1$ ranks first among
the quadratic nonresidues for $N=7$ and $N=23$, third of three for $N=8$,
third of five for $N=11$, and sixth of nine for $N=19$.  Those endpoint ranks
are snapshots, and the transition counts in
Table~\ref{tab:top-decade-spectral-fit} show why they cannot be treated as
stable asymptotic orderings.  The appropriate finite-scale conclusion is:

\begin{quote}
The unregularized races remain spectrally coherent but dynamically unsettled
through $3\times10^{14}$.  Low-zero explicit-formula truncations reproduce the
observed trajectories with high correlation, while frequent rank and leader
changes persist throughout the top decade.  These data support a low-zero
transient mechanism, not a universal stabilization threshold or an eventual
ordinary-count ordering.
\end{quote}

The distinction between statistics is essential.  Equations
\eqref{eq:ordinary-race-normalization}--\eqref{eq:ordinary-race-truncation}
describe ordinary counts.  A theorem about a separately defined
Gaussian-regularized spectral statistic would concern that statistic only.
Inferring an eventual ordering of ordinary counts would require an additional
transfer theorem with explicit error control; no such transfer is supplied by
the computation here.

\subsection{Reproducibility, data availability, and limitations}

The authoritative input is \texttt{curve\_3e14.tsv}, with SHA-256 digest
\begin{center}
\ttfamily\footnotesize
57957bdb3ce3243272c3d4b8e9ffe7dfb734b759f48b63becf7ae6f924e1caab.
\end{center}
The accompanying reproducibility package contains the PARI/GP zero generator,
the reconstruction program, fit metrics, the complete transition ledger,
per-mode attribution, figure source, a negative-test suite for the verifier,
and a SHA-256 manifest.  Its \texttt{independent\_n19} follow-up contains the
Arb sign-change certificate, the $K=25,50,100$ reconstruction, seven
adversarial verifier fixtures, and a separate manifest.  Running the two
pipeline scripts regenerates their tabular outputs and checks the declared
invariants.  For an
archival release, the package, input dataset, exact software versions, and
manifest should be deposited under a permanent DOI; the DOI replaces this
sentence before submission.

The numerical interpretation has four limits.  First, the reconstruction uses
the GRH critical-line form of the explicit formula.  Second, finite zero lists
do not prove GRH or completeness of the zero set.  Third, correlation measures
finite explanatory agreement, not mathematical implication.  Fourth, neither
the observed curves nor their truncations prove a universal or eventual
ordering.  The results are therefore best used as a controlled finite-scale
diagnostic and as a constraint on any proposed asymptotic theory.
.
% Example:
% \newcommand{\spectralfigurepath}{path/to/spectral_reconstruction.pdf}

\section{Transient spectral oscillations at the 300-trillion scale}
\label{sec:spectral-transients}

We compare the ordinary prime-count races with low-zero truncations of their
standard explicit formulas.  This calculation has two deliberately limited
purposes.  First, it tests whether low Dirichlet zeros account quantitatively
for the trajectories observed through $3\times 10^{14}$.  Second, it tests the
specific finite-scale claim that the complex-character interference has
subsided by that range.  The calculation concerns the unregularized counting
functions $\pi(x;N,a)$.  It neither proves the asymptotic proposed for the
Gaussian-regularized statistic elsewhere in this paper nor transfers an
ordering between the two statistics.

\subsection{Normalization and spectral truncation}

For a reduced residue class $a\pmod N$, set
\begin{equation}
  E_N(x;a,1)
  :=\frac{\varphi(N)\log x}{\sqrt{x}}
    \bigl(\pi(x;N,a)-\pi(x;N,1)\bigr).
  \label{eq:ordinary-race-normalization}
\end{equation}
Let
\[
  C_N(a):=\#\{b\bmod N:b^2\equiv a\pmod N\}.
\]
Assuming the critical-line form used in the standard GRH explicit formula,
we form the finite truncation
\begin{equation}
  \widetilde E_{N,K}(x;a,1)
  :=C_N(1)-C_N(a)
   -\sum_{\chi\ne\chi_0}\sum_{j=1}^{K}
     2\operatorname{Re}\!\left(
       \frac{(\overline{\chi(a)}-1)e^{i\gamma_{\chi,j}\log x}}
            {\tfrac12+i\gamma_{\chi,j}}
     \right),
  \label{eq:ordinary-race-truncation}
\end{equation}
where $\gamma_{\chi,j}>0$ is the $j$th positive ordinate returned for
$L(s,\chi)$.  Thus $K$ means the first $K$ positive ordinates for
\emph{each} nonprincipal character, rather than the first $K$ modes globally.
The coefficient $\overline{\chi(a)}-1$ is consistent with selecting the
ordered difference between the classes $a$ and $1$.  It should not be
confused with the separate character kernel used to define the regularized
statistic.

The ordinates were computed for every nonprincipal character modulo
\[
  N\in\{7,8,11,19,23\}
\]
through height $80$.  Two runs of \texttt{lfunzeros}, with mesh parameters
$64$ and $96$, were required to agree character by character.  Every returned
ordinate used in the reconstruction was also required to satisfy
\[
  \left|L\!\left(\tfrac12+i\gamma_{\chi,j},\chi\right)\right|<10^{-28}.
\]
The two mesh computations are robustness checks within one implementation,
not an independent proof that the zero list is complete.  As a separate
cross-check, all three first-zero anchors modulo $8$ agree with a pre-existing
high-precision \texttt{mpmath} computation to $10^{-12}$.

For the exceptional modulus-$19$ mode discussed below, a second implementation
constructs the character directly with Python FLINT and evaluates its Hardy
$Z$-function using Arb ball arithmetic.  Opposite, sign-definite endpoint
balls give an independent existence certificate for a critical-line zero.
This certificate does not assert uniqueness of that zero, completeness below
it, or GRH.

\subsection{Finite computational result}

The observed data consist of 438 logarithmically spaced checkpoints through
$3\times10^{14}$.  The ``top decade'' below is the 53-point subwindow
$x\ge 3\times10^{13}$.

\begin{compresult}[verified finite computation]
\label{prop:finite-spectral-computation}
For the immutable input dataset identified in the reproducibility statement
below, the checked pipeline produces 13,578 reconstruction rows.  At
$K=25$, the correlations between
$\widetilde E_{N,25}(x;-1,1)$ and $E_N(x;-1,1)$ on the 53-point top-decade
window are respectively
\[
  0.964909993,\quad 0.930923926,\quad 0.938516494,
  \quad 0.971455692,\quad 0.826435845
\]
for $N=7,8,11,19,23$.  Across the 52 consecutive top-decade intervals,
the observed rank of $-1$ changes respectively
\[
  17,\quad 23,\quad 30,\quad 17,\quad 39
\]
times.  In particular, for every one of these five moduli the class $-1$
both attains rank one and leaves rank one within the top decade.
\end{compresult}

This is a computational proposition: it records the output of a specified,
hash-checked finite calculation.  It is not a theorem about omitted zeros,
GRH, or limiting prime-race behavior.

\begin{table}[tb]
  \centering
  \small
  \caption{Top-decade reconstruction of the $-1$ race.  The correlation and
  RMSE use $K=25$ zeros per nonprincipal character.  Rank and leader changes
  refer to the observed ordinary-count curves over 52 consecutive intervals.}
  \label{tab:top-decade-spectral-fit}
  \begin{tabular}{@{}rrrrrr@{}}
    \toprule
    $N$ & correlation & RMSE & rank changes & leader changes
        & fraction $-1$ leads\\
    \midrule
     7 & 0.965 & 0.515 & 17 & 13 & 0.358\\
     8 & 0.931 & 0.626 & 23 & 15 & 0.358\\
    11 & 0.939 & 0.919 & 30 & 26 & 0.377\\
    19 & 0.971 & 1.745 & 17 &  8 & 0.415\\
    23 & 0.826 & 1.728 & 39 & 23 & 0.132\\
    \bottomrule
  \end{tabular}
\end{table}

Figure~\ref{fig:spectral-reconstruction} compares the observed curves with
the $K=1,3,10,25$ truncations.  For $N=7,11,19$, even a few modes reproduce
much of the slow structure.  For $N=23$, one zero per character is inadequate,
and the fit improves only when substantially more modes are retained.  This
behavior makes a single-mode ``onset scale'' unreliable: distinct characters
and ordinates remain active simultaneously.

\begin{figure}[p]
  \centering
  \includegraphics[width=\textwidth]{\spectralfigurepath}
  \caption{Observed normalized ordinary-count races (black), compared with
  explicit-formula truncations using the first $K=1,3,10,25$ positive zeros of
  every nonprincipal character.  The panels show $N=7,8,11,19,23$ and the
  class $a=-1\pmod N$.}
  \label{fig:spectral-reconstruction}
\end{figure}

\subsection{The very low odd-character mode modulo 19}

The computation identifies a load-bearing correction to a one-mode
description of the modulus-19 race.  For the primitive character represented
by Conrey index $13$ modulo $19$, of character order $18$, the first returned
positive ordinate is
\begin{equation}
  \gamma_{19,13,1}=0.0189563990802261\ldots,
  \qquad \chi(-1)=-1.
  \label{eq:mod19-low-zero}
\end{equation}
Both mesh runs return this ordinate, and direct evaluation at the zero passes
the residual threshold above.  Independently, Arb gives a sign-change bracket
of width $5.94\times10^{-84}$ around
$0.018956399080226142994\ldots$ and an upper bound below
$3.55\times10^{-85}$ for the modulus of the midpoint $L$-value.  The
intermediate-value theorem therefore certifies existence of at least one zero
in the bracket.  The corresponding mode is active in the $-1$ race: on the
top decade its root-mean-square contribution is $6.719$, and its correlation
with the centered observed curve is $0.728$.

The conclusion is stable under a substantially deeper truncation.  Extending
all $17$ nonprincipal characters modulo $19$ to at least $100$ positive
ordinates improves the top-decade $-1$ correlation from $0.9715$ at $K=25$
to $0.9925$ at $K=100$, while the RMSE decreases from $1.745$ to $1.252$.
The $K=100$ reconstruction retains $14$ rank changes and $7$ leader changes,
with rank and leader agreement against the observed data of $90.6\%$ and
$98.1\%$, respectively.  Increasing spectral depth therefore strengthens,
rather than removes, the transient-rank conclusion.

Consequently, an estimate based on an ordinate near $1.74$ cannot be described
as using the lowest relevant complex-character zero modulo $19$.  The much
smaller ordinate in \eqref{eq:mod19-low-zero} produces a substantially slower
oscillation that remains visible near $3\times10^{14}$.  This does not by
itself settle the ultimate behavior of the race; it shows that the proposed
finite scale is not a demonstrated point beyond complex-character
interference.

\subsection{Observed rankings and what they do not imply}

At the final checkpoint $x=3\times10^{14}$, the class $-1$ ranks first among
the quadratic nonresidues for $N=7$ and $N=23$, third of three for $N=8$,
third of five for $N=11$, and sixth of nine for $N=19$.  Those endpoint ranks
are snapshots, and the transition counts in
Table~\ref{tab:top-decade-spectral-fit} show why they cannot be treated as
stable asymptotic orderings.  The appropriate finite-scale conclusion is:

\begin{quote}
The unregularized races remain spectrally coherent but dynamically unsettled
through $3\times10^{14}$.  Low-zero explicit-formula truncations reproduce the
observed trajectories with high correlation, while frequent rank and leader
changes persist throughout the top decade.  These data support a low-zero
transient mechanism, not a universal stabilization threshold or an eventual
ordinary-count ordering.
\end{quote}

The distinction between statistics is essential.  Equations
\eqref{eq:ordinary-race-normalization}--\eqref{eq:ordinary-race-truncation}
describe ordinary counts.  A theorem about a separately defined
Gaussian-regularized spectral statistic would concern that statistic only.
Inferring an eventual ordering of ordinary counts would require an additional
transfer theorem with explicit error control; no such transfer is supplied by
the computation here.

\subsection{Reproducibility, data availability, and limitations}

The authoritative input is \texttt{curve\_3e14.tsv}, with SHA-256 digest
\begin{center}
\ttfamily\footnotesize
57957bdb3ce3243272c3d4b8e9ffe7dfb734b759f48b63becf7ae6f924e1caab.
\end{center}
The accompanying reproducibility package contains the PARI/GP zero generator,
the reconstruction program, fit metrics, the complete transition ledger,
per-mode attribution, figure source, a negative-test suite for the verifier,
and a SHA-256 manifest.  Its \texttt{independent\_n19} follow-up contains the
Arb sign-change certificate, the $K=25,50,100$ reconstruction, seven
adversarial verifier fixtures, and a separate manifest.  Running the two
pipeline scripts regenerates their tabular outputs and checks the declared
invariants.  For an
archival release, the package, input dataset, exact software versions, and
manifest should be deposited under a permanent DOI; the DOI replaces this
sentence before submission.

The numerical interpretation has four limits.  First, the reconstruction uses
the GRH critical-line form of the explicit formula.  Second, finite zero lists
do not prove GRH or completeness of the zero set.  Third, correlation measures
finite explanatory agreement, not mathematical implication.  Fourth, neither
the observed curves nor their truncations prove a universal or eventual
ordering.  The results are therefore best used as a controlled finite-scale
diagnostic and as a constraint on any proposed asymptotic theory.
